\section{Conclusion}

SME set-asides are the most widespread procurement-preference instrument worldwide, but the literature has struggled to say \emph{how} they raise prices---whether through thinner bidder pools or through less aggressive bidding within surviving auctions. This paper provides direct bid-distribution evidence on that mechanism. Exploiting the March~2018 reinterpretation of Brazilian procurement law that extended Sao Paulo's SME-only regime to medical supplies---the only previously exempt product group---I show that the winning bid immediately moved ten percentage points closer to the buyer's reference price, and the standard deviation of phase-2 submitted bids shrank monotonically in the number of participating firms, with both effects persisting under conditioning on exactly two, three, or five-plus firms. A \citet{gelbach2016} decomposition attributes only six percent of the 12--13\% price increase to entry and winner composition; ninety-four percent reflects changes in bidding behavior within auctions, exactly as the auction theory of \citet{bulow1996} and \citet{krasnokutskaya2011} predicts. Identification proceeds from a single-cohort staggered DiD with 76 never-treated controls, robust to the~\citet{callaway2021}, \citet{sun2021}, and \citet{goodmanbacon2021} estimators; the Goodman-Bacon decomposition confirms that 100\% of the coefficient comes from the clean treated-vs-never-treated 2$\times$2 comparison, and a formal test bounds the spillover bias below 1\% of the headline.

The mechanism result has two direct implications for policy design. Because 92\% of Group-65 procurement value is concentrated in the top value quartile, a simple reference-value exemption---applying the SME-only rule only to items below approximately R\$7{,}600---recovers 90\% of the fiscal cost while preserving SME access to 75\% of items by count; this is a concrete lever for the implementation of Law~14{,}133/2021, Brazil's new procurement framework, and illustrates how administrative procurement data can calibrate second-best reforms that preserve the equity goals of SME policies while curbing their fiscal footprint. Separately, CMED-regulated pharmaceuticals (class 6531) bear 62\% larger price distortion than non-pharma medical supplies, consistent with contract prices lying below the CMED cap and the SME restriction compressing competition within the residual slack. Two Brazilian regulatory regimes---SME procurement preferences and drug price controls---interact perversely: tightening either without accounting for the other leaves procurement cost higher than either instrument alone would deliver, a finding new to the literature on procurement preferences that is particularly relevant for countries such as Brazil, India, and Mexico that combine SME preferences with pharmaceutical price controls.

The aggregate fiscal cost for the SME-only rule in Group~65 alone is R\$50--85 million over the 18-month window, 7--12\% of the group's procurement value. Because BEC processes approximately R\$13 billion of procurement per year and Brazilian public procurement exceeds R\$1.5 trillion annually (approximately 10\% of GDP), proportional extrapolation across the procurement categories subject to SME-only rules indicates fiscal costs on the order of several billion reais per year at the consolidated federal-state-municipal level. The 5--11 km increase in winner-to-buyer distance under open tenders confirms that the SME preference does deliver geographic proximity, and the value-threshold reform proposed above would largely preserve that benefit.

The design has the standard limitations of a single-cohort event study and of a policy instrument that is strictly local in application. Statistical inference is limited by the single-unit treatment structure (randomization-inference $p=0.317$), so the evidence for causality rests on four-outcome consistency, the convergence of multiple estimators, and the legalistic nature of the 2018 reinterpretation rather than on a classical permutation test; external validity to other product groups, states, and procurement regimes requires the untestable assumption of comparable treatment effects. The heterogeneity documented across item value, firm age, and industry sector counsels against naive extrapolation of the \emph{magnitudes}, even as the mechanism---the contraction of competitive bidding pressure under set-asides---should generalize. Extending the RAIS window to 2018--2022 would permit evaluation of the supply-side benefits of SME preferences (employment, wages, survival) that the present paper uses only for identification robustness; structural estimation of an auction model with endogenous SME participation would quantify welfare consequences of alternative preference designs, including variants of the value-threshold reform the data here support.
